% Каркас магістерської дисертації (КПІ, ННІ ПСА, кафедра ШІ, 2026).
% Збирання: latexmk -xelatex main.tex   (або lualatex). Поруч має лежати kpi-thesis.sty.
\documentclass[14pt,a4paper,oneside]{extreport}
\usepackage{kpi-thesis}

\begin{document}

% Титул, декларація, завдання — з Word-шаблону кафедри, експортовані в PDF (F-14.1, F-4.2).
% Вони враховуються в нумерації, але номер на них не друкується.
\IfFileExists{front.pdf}{\kpifrontmatter{front.pdf}}{}

% ---------- РЕФЕРАТ / ABSTRACT (F-15.*) ----------
\structuralnotoc{Реферат}
Магістерська дисертація: 98~с., 25~рис., 12~табл., 45~посилань, 2~додатки.

РОЗПІЗНАВАННЯ МОВЛЕННЯ, НАПІВКЕРОВАНЕ НАВЧАННЯ, ПСЕВДОРОЗМІТКА, WHISPER, УКРАЇНСЬКА МОВА.

Об'єкт дослідження~-- \dots

Предмет дослідження~-- \dots

Мета роботи~-- \dots

\structuralnotoc{Abstract}
\begin{otherlanguage}{english}
Master's thesis: 98~p., 25~figures, 12~tables, 45~references, 2~appendices.

SPEECH RECOGNITION, SEMI-SUPERVISED LEARNING, PSEUDO-LABELING.

The object of the study is \dots
\end{otherlanguage}

\tableofcontents

% ---------- ПЕРЕЛІК СКОРОЧЕНЬ (за потреби, F-2.10) ----------
\structural{Перелік умовних позначень, символів, скорочень і термінів}
\abbr{ASR}{автоматичне розпізнавання мовлення (Automatic Speech Recognition);}
\abbr{WER}{частка помилково розпізнаних слів (Word Error Rate).}

% ---------- ВСТУП ----------
\structural{Вступ}
Актуальність теми зумовлена \dots~\cite{radford2023}.

% ---------- РОЗДІЛ 1 ----------
\chapter{Огляд методів розпізнавання мовлення}
У цьому розділі розглянуто \dots{} (преамбула обов'язкова, F-3.4).

\section{Моделі розпізнавання}
Архітектуру моделі показано на рис.~\ref{fig:arch}.

\begin{figure}[htbp]
  \centering
  \includegraphics[width=0.9\textwidth]{example-image}
  \figlegend{\textit{x}~-- вхідний сигнал; \textit{y}~-- транскрипт}
  \caption{Архітектура моделі}\label{fig:arch}
\end{figure}

Результати порівняння наведено в табл.~\ref{tab:wer}.

\begin{table}[htbp]
  \caption{Порівняння моделей}\label{tab:wer}
  \centering
  \begin{tabular}{|L{8cm}|C{4cm}|}
    \hline
    \textbf{Модель} & \textbf{WER, \%} \\ \hline
    Whisper large-v3 & 12,3 \\ \hline
    wav2vec~2.0 & 14,1 \\ \hline
  \end{tabular}
\end{table}

Метрику WER обчислюють за формулою~\eqref{eq:wer}:
\begin{equation}\label{eq:wer}
  \mathrm{WER} = \frac{S + D + I}{N},
\end{equation}
\begin{where}
  $S$~-- кількість замін;\\
  $D$~-- кількість видалень;\\
  $I$~-- кількість вставок;\\
  $N$~-- кількість слів еталонного транскрипту.
\end{where}

\section{Постановка задачі}
Метою роботи є \dots{} Для цього розв'язано такі завдання:
\begin{itemize}
  \item зібрати корпус нерозміченого мовлення;
  \item побудувати псевдорозмітку.
\end{itemize}

Етапи експерименту описано нижче.
\begin{enumerate}
  \item Попереднє навчання моделі.
  \item Донавчання на псевдорозмітці.
\end{enumerate}

\chapterconclusions
У розділі проаналізовано \dots

% ---------- ВИСНОВКИ ----------
\structural{Висновки}
\begin{enumerate}
  \item Проаналізовано \dots
\end{enumerate}

% ---------- ПЕРЕЛІК ПОСИЛАНЬ (ДСТУ 8302:2015, порядок першої згадки) ----------
\begin{thebibliography}{99}
  \bibitem{radford2023} Radford A., Kim J. W., Xu T. et al. Robust Speech Recognition via Large-Scale Weak Supervision. \textit{Proc. 40th International Conference on Machine Learning}. 2023. P.~28492--28518. URL: https://proceedings.mlr.press/v202/radford23a.html (дата звернення: 24.11.2026).
\end{thebibliography}

% ---------- ДОДАТКИ ----------
\kpiappendix
\chapter{Лістинг програми}
Код наведено в лістингу~\ref{lst:main}.
\begin{lstlisting}[caption={Основний модуль},label=lst:main,language=Python]
print("hello")
\end{lstlisting}

\end{document}
